% Auto-generated by extract_cwe_cve_examples.py
% Copy/paste into chap3.tex where appropriate.

\paragraph{CWE-863（Incorrect Authorization）}
\noindent\textbf{CVE-2019-11247}（Kubernetes）：The Kubernetes kube-apiserver mistakenly allows access to a cluster-scoped custom resource if the request is made as if the resource were namespaced. Authorizations for the resource accessed in this manner are enforced u...（成因：安全逻辑缺陷/2.1 访问控制缺陷）\
\noindent\textbf{CVE-2021-43858}（MinIO）：MinIO is a Kubernetes native application for cloud storage. Prior to version `RELEASE.202  2021-12-27T07-23-18Z`, a malicious client can hand-craft an HTTP API call that allows for updating policy for a user and gaining ...（成因：安全逻辑缺陷/2.3 授权验证不完善）\

\paragraph{CWE-269（Improper Access Control (Generic)）}
\noindent\textbf{CVE-2021-43858}（MinIO）：MinIO is a Kubernetes native application for cloud storage. Prior to version `RELEASE.202  2021-12-27T07-23-18Z`, a malicious client can hand-craft an HTTP API call that allows for updating policy for a user and gaining ...（成因：安全逻辑缺陷/2.3 授权验证不完善）\
\noindent\textbf{CVE-2022-29179}（Cilium）：Cilium is open source software for providing and securing network connectivity and loadbalancing between application workloads. Prior to versions 1.9.16, 1.10.11, and 1.11.15, if an attacker is able to perform a containe...（成因：配置错误/1.1 权限配置过宽）\

\paragraph{CWE-200（Exposure of Sensitive Information to an Unauthorized Actor）}
\noindent\textbf{CVE-2018-5256}（CoreOS Tectonic）：CoreOS Tectonic 1.7.x before 1.7.9-tectonic.4 and 1.8.x before 1.8.4-tectonic.3 mounts a direct proxy to the kubernetes cluster at /api/kubernetes/ which is accessible without authentication to Tectonic.（成因：配置错误/1.3 默认不安全配置）\
\noindent\textbf{CVE-2019-3869}（Tower）：When running Tower before 3.4.3 on OpenShift or Kubernetes, application credentials are exposed to playbook job runs via environment variables. A malicious user with the ability to write playbooks could use this to gain ...（成因：敏感信息泄露/3.3 配置暴露敏感信息）\

\paragraph{CWE-20（Improper Input Validation）}
\noindent\textbf{CVE-2019-11247}（Kubernetes）：The Kubernetes kube-apiserver mistakenly allows access to a cluster-scoped custom resource if the request is made as if the resource were namespaced. Authorizations for the resource accessed in this manner are enforced u...（成因：安全逻辑缺陷/2.1 访问控制缺陷）\
\noindent\textbf{CVE-2023-3676}（Kubernetes Windows）：A security issue was discovered in Kubernetes where a user that can create pods on Windows nodes may be able to escalate to admin privileges on those nodes. Kubernetes clusters are only affected if they include Windows n...（成因：安全逻辑缺陷/2.3 授权验证不完善）\

\paragraph{CWE-22（Improper Limitation of a Pathname to a Restricted Directory）}
\noindent\textbf{CVE-2022-24730}（Argo CD）：Argo CD is a declarative, GitOps continuous delivery tool for Kubernetes. Argo CD starting with version 1.3.0 but before versions 2.1.11, 2.2.6, and 2.3.0 is vulnerable to a path traversal bug, compounded by an improper ...（成因：代码注入/4.2 路径遍历）\
\noindent\textbf{CVE-2022-24877}（Flux）：Flux is an open and extensible continuous delivery solution for Kubernetes. Path Traversal in the kustomize-controller via a malicious `kustomization.yaml` allows an attacker to expose sensitive data from the controller'...（成因：代码注入/4.2 路径遍历）\

\paragraph{CWE-287（Improper Authentication）}
\noindent\textbf{CVE-2022-23652}（capsule-proxy）：capsule-proxy is a reverse proxy for Capsule Operator which provides multi-tenancy in Kubernetes. In versions prior to 0.2.1 an attacker with a proper authentication mechanism may use a malicious `Connection` header to s...（成因：安全逻辑缺陷/2.4 输入验证不足）\
\noindent\textbf{CVE-2022-29165}（Argo CD）：Argo CD is a declarative, GitOps continuous delivery tool for Kubernetes. A critical vulnerability has been discovered in Argo CD starting with version 1.4.0 and prior to versions 2.1.15, 2.2.9, and 2.3.4 which would all...（成因：安全逻辑缺陷/2.2 身份验证绕过）\

